\section{贯通算例：同一量测下的 NJ、筛选与必要性}
\label{sec:example}
取 $n=T=3,Q=0$：
\begin{equation}
 \begin{gathered}
 P=\diag(1,10,1),\quad
 Y=\begin{bmatrix}10&20&3\\2&100&1\\3&10&10\end{bmatrix},\\
 R_\star=\begin{bmatrix}10&2&1\\2&10&1\\1&1&10\end{bmatrix}
 =\diag(8,8,9)+J_{123}+J_{12}.
 \end{gathered}
 \label{eq:example}
\end{equation}
误差只有 $E_{13}=E_{31}=2$。统一缩小响应和量测可放入正常平方电压降范围，比较值相应缩放。误差位置仅用于解释构造，不提供给算法。下文有限盒例取共同上界 20，保证所有构造解均在域内；这不是生产默认参数的运行结果。

\subsection{输入权重使两个准则分离}
无约束最小二乘给出
\[
 \widehat R=YP^{-1}=\begin{bmatrix}10&2&3\\2&10&1\\3&1&10\end{bmatrix}.
\]
它对称、非负、有序，却不满足三点树关系。根锚定距离满足 $d_{0i}=10$、$d_{12}=16,d_{13}=14,d_{23}=18$。四点和为 $26,24,28$，NJ 选择错误的 $13\mid02$。

直接 L1 拟合则为
\begin{equation}
 \begin{aligned}
 L(R)={}&|R_{11}-10|+10|R_{22}-10|+|R_{33}-10|\\
 &+11|R_{12}-2|+2|R_{13}-3|+11|R_{23}-1|.
 \end{aligned}
 \label{eq:exampleloss}
\end{equation}
三个非对角权重 $11,2,11$ 来自同一批量测的输入幅值。

\begin{table}[htbp]
\centering
\begin{tabular}{lcc}
\toprule
完整结构类型 & 最小总 L1 & 最小平方损失\\
\midrule
共同块加 $\{12\}$ & $4$ & $808/103$\\
共同块加 $\{13\}$ & $11$ & $101/2$\\
共同块加 $\{23\}$ 或星形 & $13$ & $3737/68$\\
\bottomrule
\end{tabular}
\caption{在三种分支与星形的完整结构域上比较；对角均可匹配。}
\label{tab:examplecost}
\end{table}


$\{12\}$ 类型要求 $R_{13}=R_{23}=c\le R_{12}$。L1 取 $c=1,R_{12}=2$ 唯一达到 4，因为共享值 1 的权重 11 大于共享值 3 的权重 2。$\{13\}$ 类型要求 $R_{12}=R_{23}=c$，由 $11|c-2|+11|c-1|\ge11$ 得下界，并可达到。

对 $\{23\}$ 类型，写 $R_{12}=R_{13}=c$、$R_{23}=a\ge c$。用 $|a-1|\ge a-1\ge c-1$，分别在 $c\le2$ 和 $c\ge2$ 上展开另外两个绝对值的线性下界，得到
\[
 11|c-2|+2|c-3|+11|a-1|
 \ge\begin{cases}17-2c\ge13,&c\le2,\\20c-27\ge13,&c\ge2.\end{cases}
\]
$a=c=2$ 达到 13，故最优解落在星形边界。各类型的独立终端权重均可非负，因此上述下界可行。

平方损失的非对角权重为 $101,2,101$。$\{12\}$ 类型取两项加权均值 $c=(2\cdot3+101\cdot1)/103=107/103$，高值 $R_{12}=2$，得 $808/103$；$\{13\}$ 类型的共享值为 $3/2$，得 $101/2$。对 $\{23\}$ 类型，无约束最优为 $c=208/103,a=1$，违反 $a\ge c$；严格凸性使约束最优落在 $a=c$，加权均值为 $103/68$，得 $3737/68$。所以优势不依赖 L1 独有的鲁棒性，而在于共同树约束与量测权重的联合使用。

若每轮精确新增一个块，以三个 singleton 起步，L1 首轮 $\{123\},\{12\},\{13\},\{23\}$ 的目标为 $13,17,33,28$，分别由单列拟合或共同值的加权中位数取得。先选全集，再选 $\{12\}$ 达到 4，因此不依赖更全面的全树搜索。

改取 $P=\diag(1,s,1)$ 且 $Y=\widehat RP$，对 $s>1$，L1 两个主要类型的最优值为 $4$ 与 $s+1$，星形为 $s+3$，故严格选真分支当且仅当 $s>3$。平方损失令 $W=s^2+1$，两个主要值为 $8W/(W+2)$ 与 $W/2$，其差
\[
 \frac W2-\frac{8W}{W+2}=\frac{W(W-14)}{2(W+2)}
\]
给出阈值 $s>\sqrt{13}$；另一个错误类型仍落在星形边界，其值不小于错误 $\{13\}$ 类型。$s=1$ 时，两种损失都选错。若 NJ 前先执行同一量测加权树拟合，差异也会消失。

\subsection{当前最优对偶与旧证书的严格比较}
令 B 含三个 singleton 与全集。其最优核对角为 10、非对角均为 2，$J_B=13$。取
\begin{equation}
 U_0=\begin{bmatrix}0&9/10&1\\0&0&-1\\1&-1&0\end{bmatrix},\qquad
 U_1=\begin{bmatrix}0&1&0\\1&0&-1\\0&-1&0\end{bmatrix}.
 \label{eq:duals}
\end{equation}
两者均盒可行，对 B 的列内积均为零，故 $D_B(u_0)=13,D_B(u_1)=11$。候选分数为
\begin{equation}
 \begin{array}{c|rrr}
 &a_{12}^\top u&a_{13}^\top u&a_{23}^\top u\\\hline
 u_0&9&2&-11\\
 u_1&11&0&-11
 \end{array}
 \label{eq:scores}
\end{equation}
上界为 20 时，$u_0$ 的筛选界与零取最大后为 $0,0,13$，$u_1$ 则给 $\{13\}$ 下界 11。得到损失 4 的合法候选后，证书库可排除另外两个候选。

旧最优核的非零残差迫使所有旧最优对偶满足 $U_{13}=U_{31}=1$，正且未触界的 singleton 又迫使对角乘子为零。因此任一旧最优对偶都给 $a_{13}^\top u=2$，仅在当前最优对偶面挑选也无法改善这个界。$U_1$ 目标稍差却更有用，说明应回收不同已求解 LP 的证书。本文手工构造它以证明这一点，未声称取得它无需计算。

\subsection{同一个对偶得到完整禁块值}
目标 $C=\{12\}$ 的所有非目标原子只有三个 singleton、全集、$\{13\}$ 和 $\{23\}$。式~\eqref{eq:scores} 和旧列的零内积表明 $u_1$ 满足式~\eqref{eq:globaldual}，故 $J_{\neg C}\ge11$。合法 $\{13\}$ 树达到 11，因而
\begin{equation}
 J_{\neg\{12\}}=11,\qquad J_{\rm ref}=4.
 \label{eq:necessityexample}
\end{equation}
对 $4\le\varepsilon<11$，该块在完整非负共同树模型中必要；这份证书不需要有限权重上界。若用平均损失，两个值为 $4/9,11/9$，预算也应除以 9。

\section{合理可拓展性与验收边界}
\label{sec:extension}
\subsection{中间问题的作用域}
可行整数解可作同域上界；固定结构 LP 原始解可作重估起点；按式~\eqref{eq:dual} 重算的对偶可跨候选提供下界；结构提升可用于定价或全原子证书检查。当前 \texttt{dual\_bound} 只是 MIP 标量界，不是逐量测对偶 u；需导出并校验连续 LP 证书后才能实现上述复用。

若完整禁块域被区域 $\calM_q$ 覆盖，且每个区域有下界 $L_q$，则
\begin{equation}
 J_{\neg C}\ge\min_qL_q.
 \label{eq:coverage}
\end{equation}
同一完整域上的独立下界则取最大。证书至少保存数据及归一化、根/终端标签、参数界、物理先验、结构区域、合法上界、有效下界和停止原因。不同前向旧族的节点界不能不经覆盖证明就合成全树下界。

当前上界为 4、节点下界为 6 时，求最优可以剪枝；若预算为 10，该节点仍可能有可解释模型。必要性审查须保留全部下界不超过预算且未被其他证据排除的区域。逐步补入共同三点关系时，每一步仍是完整域的外松弛，下界有效；候选上界须通过全部树关系和原参数界检查。

\subsection{主要代价和失败条件}
\textbf{筛选收益取决于界的质量和复用成本。} 大上界会放大筛选惩罚，使下界退化为零；证书库和定价也有开销。加速应比较同数据、同参数域和同停止精度下的总时间、调用数与有效 gap，而非只报告筛除数量。原完整 LP 的逻辑推论再加一遍，一般不会提高理论松弛界。

\textbf{完整审查大于单块扩展。} 共同三点描述没有消除整数复杂度，全集参考模型具有指数列数。当前结果支持先做小规模参照和便宜下界，再对未决块进行完整审查，尚无大规模可用性结论。

\textbf{弱块可能确实不必要。} 取 $P=I,Q=0$，
\[
 R_a=I+3J_{123}+aJ_{12},\qquad
 R'_0=I+3J_{123}+a\diag(1,1,0),\quad 0<a<1.
\]
后者不含 $\{12\}$，且 $\sav{R_a-R'_0}=2a$。故无噪声参考损失为零时，$J_{\neg\{12\}}\le2a$；固定正预算下，足够弱的块无法确认为必要。这是允许的替代解释，不应以正权重下限排除。

\textbf{硬初始块可能排除真结构。} 锁定 $\{13\}$ 会使 $\{12\}$ 永远不可新增。必要性审查只能锁定明确物理先验或已在同域、同预算下确认的块。已确认块可加入块外响应相等等式，因为不移除预算内合法解释；数据、预算或参数域改变后须重核证明链。

\subsection{最小拓展顺序与当前结论}
第一步保留原前向求解器，为固定结构 LP 增加可核查对偶、原始/对偶目标和作用域标识，先旁路筛选而不改变候选合法域。第二步建立独立必要性审查接口，先组合全列对偶证书和多锚点 LP，再对未决块投入完整禁块审查，始终保存合法替代上界。这服务于接受、暂缓或撤销块，而非事后概率打分。

最低验收包括：三终端精确参照、四终端交叉松弛反例、正内积与秩不变边界，以及有界/无界、R/X 不相容、弱块和错误初始块。完整小规模参照中，每个下界不得超过对应最优值，每次排除不得删掉需要保留的并列结构。通过后再评价效率。

\begin{table}[htbp]
\centering\small
\begin{tabularx}{\linewidth}{@{}p{0.2\linewidth}XX@{}}
\toprule
内容 & 本稿已得到 & 尚未得到\\
\midrule
对偶筛选 & 对偶、允许重估的界、证书库和反例 & 实现、筛选频率及实际加速\\
完整域必要性 & 精确参考模型、共同三点与禁块条件、参数域对齐 & 大规模求解和松弛紧度实验\\
两方向连接 & 禁块值 11 的证书；四终端松弛反例 & 一般网络中的获取成本和覆盖率\\
\bottomrule
\end{tabularx}
\caption{数学结果、现有实现和待验证拓展分开报告。}
\end{table}


\appendix
\section{源码核对与交付复核范围}
\label{app:source}
核对日期为 2026-10-28，基准提交为 \texttt{10b17c16}。工作区另有既有修改；本稿只读取相关模型文件，不声称整个工作区等于该提交。所用源文件 SHA-256 保存在 \path{verification/source_snapshot.json}。

\begin{table}[htbp]
\centering\small
\begin{tabularx}{\linewidth}{@{}p{0.42\linewidth}X@{}}
\toprule
文件或位置 & 本稿使用的事实\\
\midrule
\texttt{preprocessing.py}:8--25 & 观测根平方电压差，根表误差仍进入目标\\
\texttt{laminar\_l1\_milp.py}:574--591 & 固定块的 P/Q 特征\\
同文件 712--744 & 权重盒、平均 L1 目标、固定族 LP\\
同文件 747--810 & 乘积提升、层状关系和去重\\
同文件 1187--1198 & 初始块保留、单块扩展、验证回选\\
同文件 1329--1375 & 单步界的适用域；触界扩大并重启\\
同文件 40--80、887--905 & 返回权重、目标与 MIP 标量诊断，未导出逐行对偶库\\
\texttt{pipeline.py}:332--337 & 初始族和共同权重上界传入扩展器\\
\bottomrule
\end{tabularx}
\end{table}

除 \texttt{pipeline.py} 外，表中文件位于 \path{rnj_wzzt_core/rnj_wzzt/estimation/}；\texttt{pipeline.py} 位于包根目录。生产基线、AC 实验及项目优化程序均未为本稿运行。

本目录的 \path{verify_derivations.py} 仅用有理数复核构造解、目标、对偶内积、距离和及反例恒等式，不调用优化求解器。它不替代一般定理证明，也不是算法性能实验。引文依据是 \citet{bravo2009} 的 \S2 和 Table~1、\citet{fercoq2015} 的 \S1--3，以及 \citet{mihaescu2007} 的 \S1--3，仅使用实际阅读部分，不声称综述穷尽相关工作。
